\chapter{Event Organizer Rules}

\section{Venue}

	\begin{enumerate}[leftmargin=1.5em,labelsep=*]
		\item Depending on the Road Racing discipline, courses are usually held on one of the following options:
		\begin{enumerate}[label*=\arabic*,leftmargin=1.5em,labelsep=*]
			\item Fixed and Free Distance Races are usually held on roadways or bike paths.
			It is recommended to use a point-to-point courses without loops.
			For Fixed Distance Races additional requirements are described in \ref{sec_venue_fixed_distance}.
			\item Criteriums are held around city block(s) or within a large parking lot, allowing the course to have left 
      and right turns and multiple laps with a recommended lap length of 500 m to 1000 m.
			\item Distance-based Time Trials in officiall competitions are held on the same courses as Fixed and Free Distance Races.
			Outside of officiall competitions they can also held on the same courses as Time-based Time Trials.
			\item Time-based Time Trials must take place on a closed circuit, either on a road or a track with a recommended lap length of 200 m to 1000 m.
		\end{enumerate}
		\item The course must be designed with rider safety as a top priority.
		It should include clear signage, road barriers where necessary, and an adequate number of marshals to guide and protect participants.
		\item The road surface should be smooth, durable, and free of hazards such as deep potholes, cracks, off-camber sections, or uneven terrain.
		Rough surfaces, narrow segments, and other obstacles should be avoided or adequately managed to ensure rider safety.
		If bicycle paths are used, they must be sufficiently wide to accommodate riders at speed and should be free of any hazards.
		Objects adjacent to the course—such as benches, trees, or metal posts—must be considered potential risks and properly safeguarded.
		\item Throughout most of the course, riders must have enough space to overtake safely.
		This is not nesseccary for individual Time Trials outside of official competitions.
	\end{enumerate}

\subsection{Start and Finish Area for Races with Mass Start or Wave Start}
	\begin{enumerate}[leftmargin=1.5em,labelsep=*]
		\item The start area must be wide enough to accommodate several riders launching simultaneously in safe conditions.
		\item It is recommended that the finish area is not downhill, as downhill finishes increase the risk of crashes.
		Organizers may choose to place the finish on flat or uphill terrain.
		The final few hundred meters of the course should be straight, with no sharp turns, to allow for safe sprint finishes.
	\end{enumerate}

\subsection{Fixed Distance Races \label{sec_venue_fixed_distance}}
	\begin{enumerate}[leftmargin=1.5em,labelsep=*]
		\item For any Fixed Distance Race according to 3B.2.1, it is recommended that the course complies with the requirements for IUF World Records.
		In short:
		\begin{enumerate}[label=\roman*,leftmargin=1.75em,labelsep=*]
			\item The start and finish points of the course, measured along a theoretical straight line between them, shall not be further apart than 50\% of the race distance.
			\item The overall decrease in elevation between the start and finish shall not exceed 1:1000, i.e.\ 1m per km (0.1\%).
		\end{enumerate}
		Please refer to the IUF World Records Guidelines for details.
	\end{enumerate}

\section{Officials}

The host must designate the following officials for each road race:
\begin{itemize}
\item Race Director
\item Referee
\item Starter
\end{itemize}

\section{Communication}

\begin{comment2016}
Some thoughts on additional required communication:
\begin{itemize}
\item any changes to the standard rules
\item heat assignment/seeding method
\item false start method
\item cut-off time
\item age groups
\item starting method
\item whether the event qualifies for a world record
\item timing method
\item course map
\end{itemize}
\end{comment2016}

The host must announce the false start method at least two months before the event.

Details of all non-track racing events, or other events with unique courses or details must be published as soon as they are known.
This is to provide competitors with the information they need to train, and to help them prepare the appropriate unicycles.
These are major needs for attendees from far away.
Necessary details depend on the event, but include things like course length, elevation and elevation change, steepness, 
level of terrain difficulty, amount of turns, riding surfaces, course width, etc.
Maps should be provided when possible.
While sometimes courses cannot be planned until weeks or days before the convention, as soon as they are known the details 
must be posted to the convention web site and/or all places where convention information is posted.
It is acceptable to publish tentative courses while waiting for permits to be approved, etc.

\section{Age Groups}

The following age groups are the minimum required by the IUF to be offered at the time of registration for any Road Racing discipline: 0-13, 14-18, 19-29, 30-UP.
For any discipline for which there is a 24 Class wheel size category, also an age group 0-10 (20 Class) must be offered.
All age groups must be offered as mens and womens age group.

\section{Practice}

\begin{comment2016}
The following rule came from chapter 1.
This rule is poorly written.
What happens between 1 day and 6 days?
\end{comment2016}

If the course is open for practice to all riders for at least 7 days leading up to the event, then there are no restrictions on who can compete.
If the course is not open for practice until the day of the event, then anyone who has pre-ridden the course is not allowed to compete.
Organizers must therefore ensure that course marking and set-up are done by non-competing staff/volunteers.

\section{Race Configuration}

Riders are usually divided by age group and unicycle type, such as 24 Class versus 29 Class unicycles, 
and/or Standard (a Regular Unicycle, any size wheel and cranks) versus Unlimited.

\section{Starting Order}

%TODO Change "heat" -> "wave"
\begin{enumerate}[leftmargin=1.5em,labelsep=*]
	\item The goal in determining the starting order is to sort racers fairly by speed while still making sure that men and women race amongst themselves.
	Unless otherwise noted below, the fastest riders start first, and also within a start group (heat or mass start), 
  riders should be positioned in the line-up by speed with the fastest in front.
	\item At Unicons, the Starting Order is determined based on an expected time or average speed declared by each competitor at registration. 
  This declared value may be based on a previous race, a training session, or a personal estimate. 
  It may be entered directly or selected from a predefined range provided during registration. 
  Riders with the fastest expected performances will be placed in the earliest waves for wave starts, or at the front of the start line for mass starts. 
  Riders who do not provide an expected time or speed may be placed in the last wave or at the back of the start line.
	\item If results from a previous road race held during the same competition are available—such as a 10k before a 
  marathon—they may also be used to help determine the starting order.
	\item The starting order is the joint responsibility of the Race Director and the Referee.
\end{enumerate}

\section{Starting Configuration}

Line-up order and heats must be assigned prior to the race.
There are three allowable formats for designating the starting configuration of a Road Race: 
individual start (section \ref{subsec:road_heat-assignment_individual-start}), 
heat start (section \ref{subsec:road_heat-assignment_heat-start}), or mass start (section \ref{subsec:road_heat-assignment_mass-start}).

To determine which start configuration to use, read the following rules from top to bottom.
Once you have an outcome, \emph{disregard} the remaining rules.
\begin{itemize}
\item If this is an ``Individual Time Trial'' format race, use individual start.
\item If the course is too narrow to allow for racers to safely and fairly start in heats, use individual start.
\item If you cannot safely start five or more riders across, use individual start.
\item If the starting field consists of 30 riders or less, use a mass start.
\item If the course does not allow for ten riders to ride abreast for at least 500 meters before the course narrows, use heats of 12 or more riders.
\item If the starting field consists of more than 100 riders, use heats of 20 or more riders.
\item In all other cases, use a mass start.
\end{itemize}

Unlimited racers should start first, unless there is no risk that Unlimited riders have to pass Standard riders (for example they race on different days).

In the sections below, ``fastest rider'' means ``fastest rider by seed time.'' Seed time is defined as an estimated finish time, 
preferably based on past performance in similar event(s).
If no seed time is submitted by the rider or their coach, the organization can assign a seed time.

\subsection{Individual Start \label{subsec:road_heat-assignment_individual-start}}

Each rider is individually started at a fixed time interval, such as every 20 or 30 seconds.
Riders are sorted by speed with the fastest rider going first.
(Except in the case of an Individual Time Trial, where the race can start with either the fastest or slowest rider.)

\subsection{Wave Start \label{subsec:road_heat-assignment_heat-start}}

\begin{enumerate}[leftmargin=1.5em,labelsep=*]
	\item A wave start is a start where a smaller group of competitors start the race together.
	This type of start is commonly used in fixed distance races (e.g. 10km or Marathon) when a mass start is not feasible due to space limitations.
	\item The various classes may share the race course, but Standard racers should always start separately from Unlimited racers.
	\item Waves should consist of at least 12 riders, either men or women (no mixed heats).
	Waves may vary in size.
	Waves are sorted by speed with the fastest wave going first.
	The first wave should be devoted to the fastest men, and the second wave to the fastest women.
	These waves must be single-gender.
	The top men and the top women must have equivalent racing conditions.
	The following waves should be sorted by speed.
	The time intervals between waves 1, 2, and 3 should be set up such that following waves have the least chance of interfering with the top men and women riders.
	In all cases, the minimum time interval between two waves is 1 minute.
\end{enumerate}

\subsection{Mass Start \label{subsec:road_heat-assignment_mass-start}}

\begin{enumerate}[leftmargin=1.5em,labelsep=*]
	\item A mass start is a grouped start where all participants start the race together.
	This type of start is commonly used for long-distance races (marathon or longer free distance road races).
	\item The organizer may choose to separate the starts based on class and gender, allowing for up to four separate starts, resulting in three possible configurations:
	\begin{enumerate}[label*=\arabic*,leftmargin=1.5em,labelsep=*]
		\item Standard and Unlimited start together, all genders combined (1 start) (not suitable for road races at UNICON),
		\item Standard and Unlimited start separately, all genders combined (2 starts),
		\item Standard and Unlimited start separately, with men and women also starting separately (4 starts).
	\end{enumerate}
	\item In any case the starting order on the line may be determined based on seed times or other criteria set by the organizer.
\end{enumerate}

\section{Timing, Photo Finish}

	\begin{enumerate}[leftmargin=1.5em,labelsep=*]
		\item The use of Fully Automatic Timing and Photo Finish Systems or Transponder Timing Systems that fulfill the 
    specified criteria are permitted for timekeeping in road races at IUF sanctioned events.
		For other competitions, these are highly recommended, but Hand Timing that fulfills the specified criteria is also allowed here.
		\item The use of Fully Automatic Timing and Photo Finish Systems is permitted, 
    provided they meet the criteria set out in Rule 2D.8 and are operated accordingly.
		\item The use of Transponder Timing Systems is permitted, provided that:
		\begin{enumerate}[label=\roman*,leftmargin=2.25em,labelsep=*]
			\item None of the equipment used at the start, on the course or at the finish line is a significant obstacle to the progress of a rider.
			\item The weight of the transponder and its housing carried or worn by the rider is not significant.
			\item The System is started by the start signal.
			\item The System requires no action by a rider during the competition, at the finish or at any stage in the result processing.
		\end{enumerate}
		\item The use of Hand Timing is permitted, provided that:
		\begin{enumerate}[label=\roman*,leftmargin=2.25em,labelsep=*]
			\item Timekeepers are in line with the finish and all have a good view of the finish line.
			\item Timekeepers use manually operated electronic timers with digital readouts.
			All such timing devices are termed “watches” for the purpose of the Rules.
			\item The time is taken from the start signal.
			\item Two official Timekeepers time the winner of every event and any performance for record purposes.
			\item Each Timekeeper acts independently and without showing their watch to, or discussing their time with, any other person.
			\item The time measurement must be carried out with a resolution of at least 1/10 of a second, unless the measurement system 
      ensures that the times are always given to at least the next longer full second.
			Unless the time is an exact full second, the time will be converted to the next longer full second.
			If the measuring system only displays full seconds and it cannot be ensured that this is the next longer full second, one second must be added to the displayed time.
			If two watches are prescribed and, after converting as indicated above and two watches are prescribed, the two watches disagree, the longer time will be official.
		\end{enumerate}
		\item For all times to be official, unless the measured time is an exact whole second, 
    the time shall be converted and recorded to the next longer whole second, e.g. 1:33:47.153, the official time is recorded as 1:33:48.
		All official times must be determined considdering paragraph 6 and 7 where required.
		\item The placing order must be determined at the finishing moment in accordance with rule 3B.5.8.
		It is recommended that Judges and/or video recording(s) also be provided to assist in determining the finishing order and the identification of riders.
		For all road races at IUF sanctioned events, a Photo Finish System according to paragraph 3 or a 
    video recording perpendicular to the finish line must be available to determine the placings.
		\item If a reliable deternmination of the finishing order between competitors is not possible, the competitors must be ranked equally and assigned the same official time.
		\item Only official times will be published.
		\item When two competitors have the same official time and the placing is determined based on the 
    actual finishing order, the result list should indicate if a Photo Finish System was used or not (e.g., "Photo Finish: +0.03" or "Timekeeper Decision").
	\end{enumerate}

\section{Optional Race-End Cut-Off Time}

It may be necessary to have a maximum time limit for long races, to keep events on schedule.
When this is planned in advance, it must be advertised as early as possible, so attending riders will know of the limit.
Additionally, at the discretion of the Racing Director, a race cut-off time may be set on the day of or during an event.
The purpose of this is to allow things to move on if all but a few slow racers are still on the course.
These cut-offs need not be announced in advance.
At the cut-off time, any racers who have not finished will be listed as incomplete (no time recorded, or same cut-off time recorded for all).
Optionally, if there is no more than one person on the course per age category and awards are at stake, they can be given the last place in the finishing order.
But if each participating age category has had finishers for all available awards (no awards at stake), there is no need to wait.

\section{Lead Vehicle}

\begin{enumerate}[leftmargin=1.5em,labelsep=*]
	\item A lead vehicle is used to guide the front of the race, provide visibility, and ensure safety, especially in races held on open roads.
	\item Its use is recommended for all open-road races, except time trials.
	\item The lead vehicle must stay approximately 50 to 100 meters ahead of the leading unicyclist, without interfering with the race or providing a drafting advantage.
	It must be fast and maneuverable enough to stay ahead throughout the race.
	\item A two-wheeled vehicle is preferred.
	This may be a motorcycle, an e-bike, or a fit elite cyclist capable of riding faster than the leading unicyclist.
	The use of a car is discouraged due to its width and the larger slipstream it creates.
	\item It should be clearly marked and easily identifiable by both competitors and marshals.
	\item The operator should be in contact with race officials to report hazards or incidents on the course.
	\item If no lead vehicle is feasible, alternative safety measures must be put in place.
\end{enumerate}

\section{Race Distances and Distance Measurement}

\subsection {Distance Measurement for Fixed Distance Races}

\begin{enumerate}[leftmargin=1.5em,labelsep=*]
	\item For fixed-distance races, the course must be measured accurately along a path 30 cm from the kerb or other solid 
  boundaries to the riding surface, to ensure that the measured distance is at least equal to the advertised distance.
	\begin{enumerate}[label*=\arabic*,leftmargin=1.5em,labelsep=*]
		\item When ever possible an official measurement must be carried out before the race.
		For this the course distance must be measured using a Jones Counter, following the official IAAF recommendations or an equivalent method offering comparable accuracy.
		This measurement ensures the validity of the results and the potential recognition of records.
		\item If an official measurement before the race is not feasible, a simplified method using a Jones Counter may be used.
		This method maintains a high level of accuracy and ensures a conservative measurement by applying a correction factor of 1.003.
		No records can be recognized by this measurement.
	\end{enumerate}
	\item If the simplified method was used before the race, an official measurement must be conducted after the event to validate any potential record.
\end{enumerate}

\subsubsection {Simplified method using a Jones Counter}

\begin{enumerate}[leftmargin=1.5em,labelsep=*]
	\item The following equipment is required for measurements using a Jones Counter:
	\begin{enumerate}[label=\roman*,leftmargin=1.75em,labelsep=*]
		\item A certified Jones Counter, mounted on a bicycle wheel (preferably with a smooth, well-inflated tire).
		\item A calibration course: a straight, flat, and stable section of road with a known precise length, ideally between 300 m and 500 m.
		\item A certified steel measuring tape of at least 30 meters, used to confirm or verify the exact length of the calibration course.
		\item Markers to clearly define the start and end points of the calibration course.
		\item A notebook or app to record all measurements.
	\end{enumerate}
	\item The simplified method using a Jones Counter consists of three steps, according to i - iii.
	Calibration and course measurement must be carried out consecutively at a time of day when there is no great variation in temperature.
	\begin{enumerate}[label=\roman*,leftmargin=1.75em,labelsep=*]
		\item \textbf{Calibration of the Jones Counter:} a critical step\\
		Calibrating the Jones Counter is essential to ensure the accuracy of all subsequent measurements.
		This step must be performed meticulously:
		\begin{enumerate}[label*=.\roman*,leftmargin=3.25em,labelsep=*]
			\item Setting up the calibration course:
			Choose a straight, stable section of road, and confirm its exact length using a steel measuring tape.
			Ensure that the start and end points are clearly marked and fixed.
			\item Measuring the calibration course:
			Perform 4 measurements (2 round trips) of the calibration course with the Jones Counter:
			Ride slowly and steadily along the full length of the course, maintaining a straight line.
			Record the Jones Counter readings at each end for every pass.
			\item Managing errors:
			If any reading deviates significantly from the others, it must be discarded.
			Such errors may arise from deviations in the bicycle's path or instability.
			If necessary, repeat the 4 measurements to obtain reliable data.
			\item Calculating the calibration factor (CF):
			Average the valid readings to determine the calibration factor:\\
			CF = Average number of Jones Counter counts over 4 measurements / Length of the calibration course in meters
			\item Applying the correction factor:
			Apply a correction factor of 1.003 to ensure a conservative measurement:\\
			CFcorrected = CF x 1.003
		\end{enumerate}
		\item \textbf{Measuring the course:}
		\begin{enumerate}[label*=.\roman*,leftmargin=3.25em,labelsep=*]
			\item Starting point:
			Position the bicycle equipped with the Jones Counter at the starting point of the course.
			\item Measurement:
			Ride slowly and steadily along the shortest possible path of the course (the "ideal line").
			\item Final reading:
			Record the Jones Counter reading at the finish point.
		\end{enumerate}
		\item \textbf{Calculating the measured distance:}\\
		Use the corrected calibration factor to convert the Jones Counter readings into the actual distance:\\
		Actual distance (m)= Total number of Jones Counter counts over the course / CFcorrected
	\end{enumerate}
\end{enumerate}

\section {Road Racing Events at Unicon}

\begin{enumerate}[leftmargin=1.5em,labelsep=*]
	\item It is expected that Unicon will have at least three road racing events, of which at least one is a recognized fixed distance event.
	If only one fixed distance event is held, it should be the 10km.\\
	\textbf{Note:} Examples of road racing events at Unicon include:\\
	-- a 10km,\\
	-- a Marathon or a longer free distance race with a mass start,\\
	-- a Criterium.\\
	A climbing road race may be considered instead of, or in addition to, the Criterium if the location allows for such a course.
\end{enumerate}
